% !Mode:："TeX:UTF-8"
\chapter{模板使用}
\section{模板文件结构\label{sec:files}}
整个模板根目录的文件列表如下：
\begin{center}
\begin{tabular}{|l|l|l|}
\hline
HUNNUthesis.cls & ---HUNNUthesis宏包 & \textcolor{red}{{*}}\\
\hline
hunnubib.bst & ---引文样式文件 & \textcolor{red}{{*}}\\
\hline
references/reference.bib & ---bib数据库 & \textcolor{red}{{*}}\\
\hline
figures/logo.eps & ---湖南师范大学彩色字体 & \textcolor{red}{{*}}\\
\hline
main.tex & ---\TeX{}样例文件 &\textcolor{red}{{*}}\\
\hline
\end{tabular}
\end{center}
注： \textcolor{red}{{*}} 表示\LaTeX{}模板必须的文件。
\section{示例}
对于论文中最常使用的一些功能在本节中给出示例。
\subsection{公式}
\begin{equation}
\hat{H}=\frac{\epsilon}{2}\hat{\sigma}_{z}-\frac{\Delta}{2}\hat{\sigma}_{x}+\sum_{k}\omega_{k}\hat{b}_{k}^{\dagger}\hat{b}_{k}+\sum_{k}\frac{g_{k}}{2}\hat{\sigma}_{z}(\hat{b}_{k}+\hat{b}_{k}^{\dagger})\label{eq:sbm}
\end{equation}
根据公式\ref{eq:sbm}可知，这个是对公示的引用。
\begin{align}\label{eqn3:9}
\int_{-\infty}^{+\infty}S(\tau,f)\,\mathrm{d}\tau&=\int_{-\infty}^{+\infty}x(t)\left\{\int_{-\infty}^{+\infty}\frac{|f|}{\sqrt{2\pi}}e^{\frac{-|f|^2(\tau-t)^2}{2}}\,\mathrm{d}\tau\right\}e^{-j2\pi ft}\,\mathrm{d}t\notag\\
&=\int_{-\infty}^{+\infty}x(t)e^{-j2\pi ft}\left\{\int_{-\infty}^{+\infty}\frac{1}{\sqrt{\pi}}e^{-\left[\frac{|f|(\tau-t)}{\sqrt{2}}\right]^2}\,\mathrm{d}\frac{|f|(\tau-t)}{\sqrt{2}}\right\}\,\mathrm{d}t
\end{align}
令$\theta=\frac{|f|(\tau-t)}{\sqrt{2}}$，则式\eqref{eqn3:9}可改写为
\begin{align}\label{eqn3:10}
\int_{-\infty}^{+\infty}S(\tau,f)\,\mathrm{d}\tau&=\int_{-\infty}^{+\infty}x(t)e^{-j2\pi ft}\,\mathrm{d}t\frac{1}{\sqrt{\pi}}\int_{-\infty}^{+\infty}e^{-\theta^2}\,\mathrm{d}\theta\notag\\
&=\int_{-\infty}^{+\infty}x(t)e^{-j2\pi ft}\,\mathrm{d}t\frac{2}{\sqrt{\pi}}\int_{0}^{+\infty}e^{-\theta^2}\,\mathrm{d}\theta\notag\\
&=\int_{-\infty}^{+\infty}x(t)e^{-j2\pi ft}\,\mathrm{d}t\notag\\
&=X(f)
\end{align}
\subsection{表格}
\begin{table}[H]
	\begin{center}
		\caption{希腊字母表\label{tab:Greek}}
		\begin{tabular}{|c|c|c|c|c|}
			\hline
			Alpha & Beta & Gamma & Delta & Theta\\
			\hline
			$\alpha$ & $\beta$ & $\gamma$ & $\delta$ & $\theta$\\
			\hline
			$A$ & $B$ & $\Gamma$ & $\Delta$ & $\Theta$\\
			\hline
		\end{tabular}
		\end{center}
\end{table}
这是对表\ref{tab:Greek}的引用

为了测试模型的性能，采用以下六个基准语料库对所提出的模型进行实验分析和验证。其中，表 \ref{tab:datasets} 为数据集的详细统计表。%图 \ref{fig:R8-partial-datasets}展示了部分R8数据集。

\begin{table}[htbp]
	\caption{数据集汇总}\label{tab:datasets}
	\vspace{0.5em}\centering\zihao{5}
	\renewcommand\arraystretch{1.5}
	\begin{tabular}{lcccccccc}
		\toprule[1.5pt]
		数据集 & 总文档 & 训练集 & 测试集 & 字典大小 &节点数  &类数 & 平均长度 & 稀疏度\\
		\midrule[1pt]
		R8     & 7674  & 5485  & 2189  & 7688  & 15362 & 8   & 65.72 & 234\\
		R52    & 9100  & 6532  & 2568  & 8892  & 17992 & 52  & 69.82 & 258\\
		20NG   & 18846 & 11314 & 7532  & 42757 & 61603 & 20  & 221.26  & 278\\
		Ohsumed& 7400  & 3357  & 4043  & 14157 & 21557 & 23  & 135.82 & 159\\
		MR     & 10622 & 7108  & 3554  & 18764 & 29426 & 2   & 20.39 & 1443\\
		WebKB  & 4199  & 2803  & 1396  & 6300  & 10499 & 4   & 131.97 & 80\\
		\bottomrule[1.5pt]
	\end{tabular}
	\vspace{\baselineskip}
\end{table}


\subsection{图形}
这个示例为插入图片：
\begin{figure}[H]
	\centering
		\includegraphics[width=0.618\textwidth]{logo}%图片名称，放在/figures目录下
	\caption{图片插入\label{fig:fig}}
\end{figure}
具体代码：
\begin{verbatim}%抄写环境
\begin{figure}[H]
\centering
\includegraphics[width=0.618\textwidth]{logo}%图片放在/figures目录下
\caption{图片插入\label{fig:fig}}
\end{figure}
\end{verbatim}
\begin{figure}[H]
\centering
		\includegraphics[width=0.618\textwidth]{logo}
\caption{湖南师范大新大学校名标准字\label{fig:hunnu}}
\end{figure}

\begin{figure}[htbp]
	\centering
	\subfigure[]{\includegraphics[width=0.49\linewidth]{logo} \label{fig:logo}}
	\subfigure[]{\includegraphics[width=0.49\linewidth]{logo} \label{fig:logo}}
	\subfigure[]{\includegraphics[width=0.49\linewidth]{logo} \label{fig:logo}}
	\subfigure[]{\includegraphics[width=0.49\linewidth]{logo} \label{fig:logo}}
	
	\caption{(a)为logo，(b)为logo，(c)为logo，(d)为logo}
\end{figure}


对于图\ref{fig:fig}和图\ref{fig:hunnu}的引用。
\subsection{引用}
\subsubsection{交叉引用}
对所有需要引用的公式、表格、图形，执行插入--标签后，即可使用插入-- 交叉引用自动产生引用。
\begin{itemize}
\item 哈密顿量见方程~\eqref{eq:sbm}。
\item 希腊字母表见表~\ref{tab:Greek}。引用格式与方程引用格式不同
\item 校名标准字如图~\ref{fig:hunnu}。 引用格式与方程引用格式不同
\end{itemize}
具体见代码：
\begin{verbatim}
\begin{itemize}
\item 哈密顿量见方程~\eqref{eq:sbm}。
\item 希腊字母表见表~\ref{tab:Greek}。引用格式与方程引用格式不同
\item 校名标准字如图~\ref{fig:hunnu}。 引用格式与方程引用格式不同
\end{itemize}
\end{verbatim}
\subsubsection{文献引用}
将引文的bib数据库(默认文件名为reference.bib)放入模板根目录下的references文件夹，即可通过插入--文献引用自动产生引文。
\begin{itemize}
\item Journal：An article \upcite{ELIDRISSI94,MELLINGER96,SHELL02,cnarticle}。
\item Book：An book \cite{IEEE-1363,tex,companion}。
\item Conference：A conference \cite{kocher99,DPMG,cnproceed}。
\item Manual:A manual\cite{NPB2}.
\item MasterThesis:\cite{zhubajie,metamori2004,shaheshang,FistSystem01}.
\end{itemize}
\subsection{伪代码实现}
\begin{algorithm}
\caption{放进冰箱的大象}\label{算法实例}
\begin{algorithmic}
	\REQUIRE 有一只大象
	\ENSURE 放进冰箱里
	\FOR {没有剩余的大象}
	\IF {大象比冰箱大}
	\STATE 把大象分割
	\ENDIF
	\ENDFOR
	\STATE 第一步
	\STATE 第二步
	\STATE 第三步
\end{algorithmic}
AAA\end{algorithm}
\subsection{代码展示}
可以把你的程序添加到附录里，展示自己的工作。
\begin{lstlisting}[language={[ANSI]C}, numbers=left]
#include <stdio.h>
int main(int argc, char ** argv)
{
/*打印Hello,world*/
printf("Hello, world!\n");

return 0;
}
\end{lstlisting}
\section{依赖}
HUNNUthesis依赖于以下宏包，这些宏包在常见的\LaTeX{}发行版中都包括，在安装使用之前，请确定你的\TeX{}发行版中都已正常安装这些宏包
\begin{table}[H]
	\centering
	\begin{tabular}{cccc}
		\hline
		{footmisc} &  {amsmath} &  {amsfonts} &  {amssymb} \\
		
		{graphicx} &  {svgnames} &  {xcolor} &  {mathptmx} \\
		
		{float} &  {fontenc} &  {fancyhdr} &  {lastpage} \\
		
		{etoolbox} &  {fancy} &  {caption} &  {array} \\
		
		{makecell} &  {forloop} &  {xstring} &  {hyperref} \\
		
		{tabularx} &  {enumitem} &  {ntheorem} &  {algorithm}\\
		
		{algorithmic} &  {bibentry} &  {xeCJK} &  {CJK} \\
		{listings} &  {courier} &  {} &  {} \\
		\hline
	\end{tabular}
\end{table}
如果你尚未安装这些宏包，可以启动你的 \TeX{} 发行版的宏包管理器
来安装；或者到 \url{http://www.ctan.org} 上搜索下载并安装。
\section{基本设置}
\begin{enumerate}
	\item 图片搜索路径默认设置为模板根目录下的figures/。
	\item bib数据库默认设置为模板根目录下的references/reference.bib。 其中bib文件可由任意文献库管理软件自动生成
\end{enumerate}


